\NativeTitle{33}{第五单元素养评价（A）}{（多边形的认识）}{90}{6}{一 & 二 & 三 & 四 & 五 & 六}{ &  &  &  &  & }
\NativeSlot{283}{36}{%
\ExamGroup{一、填空。{\normalfont\kaishu\fontsize{10.8bp}{14bp}\selectfont （每空2分，共38分）}}
% p033_s01
% p033_s02
\ExamQuestion{1.}{5}%
一个三角形中最少有\mbox{（\hspace{1.8em}）}个锐角，最多有\mbox{（\hspace{1.8em}）}个钝角。\par
}
\NativeSlot{333}{36}{%

% p033_s03
% p033_s04
\ExamQuestion{2.}{5}%
三角形的内角和是\mbox{（\hspace{1.8em}）}，任意两边长度的和\mbox{（\hspace{1.8em}）}第三边。\par
}
\NativeSlot{359}{36}{%
% p033_s05

\ExamQuestion{3.}{5}%
有一条腰垂直于梯形的上下两底，这样的梯形叫作\mbox{（\hspace{1.8em}）}。\par
}
\NativeSlot{384}{36}{%
% p033_s06

\ExamQuestion{4.}{5}%
在一个长方形中画一条线段，将长方形分成两个梯形，每个梯形都是\mbox{（\hspace{1.8em}）}梯形。\par
}
\NativeSlot{435}{72}{%
\par\noindent%
\ExamTile{.485}{%
% p033_s07

\ExamQuestion{5.}{5}%
下图中，三角形的高所对应的底边长是\mbox{（\hspace{1.8em}）}厘米。\par
\BookFigure{figures/p033_s07-01.tex}
}\hfill%
\ExamTile{.485}{%
% p033_s08


下图中，平行四边形\(AB\)边对应的高是\mbox{（\hspace{1.8em}）}厘米。\par
\BookFigure{figures/p033_s08-01.tex}
}\par

}
\NativeSlot{563}{45}{%
% p033_s09

\ExamQuestion{6.}{5}%
一个等腰三角形的顶角是\(60^\circ\)，它的一个底角是\mbox{（\hspace{1.8em}）}°，按角分，这个三角形也是\mbox{（\hspace{1.8em}）}三角形。\par

% p033_s10

\noindent 一个直角三角形的一个锐角是\(45^\circ\)，另一个锐角是\mbox{（\hspace{1.8em}）}°，这个三角形也是\mbox{（\hspace{1.8em}）}三角形。\par
}
\NativeSlot{640.0}{44}{%
% p033_s11

\ExamQuestion{7.}{5}%
AI机器人编程实践课上，工程师给智能机器人平行四边形感应探测底座设定尺寸，它相邻的两条边长分别是5 cm和3 cm，这个探测底座的周长是\mbox{（\hspace{1.8em}）}cm。\par
}
\NativeSlot{717}{36}{%
% p033_s12

\ExamQuestion{8.}{5}%
创客社团制作等腰三角形磁吸书签，书签的两条边长分别是4厘米和9厘米，这个书签的周长是\mbox{（\hspace{1.8em}）}厘米。\par
}
\NativeSlot{768}{36}{%
% p033_s13

\ExamQuestion{9.}{5}%
在一个三角形中，如果有两个锐角的度数和大于\(90^\circ\)，那么这个三角形一定是\mbox{（\hspace{1.8em}）}三角形。\par
}
